概率题最常见的失败发生在动笔之前。两个人拿同一份数据、同一句题面，各自算得毫无破绽，报出的数却对不上——问题不在算式，在他们心里那个``一次试验产出什么''从一开始就不是同一个东西。

\begin{center}
\begin{tikzpicture}[x=1cm,y=1cm,font=\footnotesize,>=stealth]
  \draw[black!45,fill=black!5] (-2.9,0) rectangle (2.9,0.95);
  \node at (0,0.475) {``随机挑一次访问，看有没有下单''};
  \draw[->,black!60] (-1.4,-0.05) -- (-3.0,-1.0);
  \draw[->,black!60] (1.4,-0.05) -- (3.0,-1.0);
  \draw[black!50,fill=blue!10] (-5.4,-3.1) rectangle (-1.0,-1.2);
  \draw[black!50,fill=blue!10] (1.0,-3.1) rectangle (5.4,-1.2);
  \node[black!75] at (-3.2,-1.55) {甲：一条记录 \(=\) 一个用户};
  \node[black!75] at (3.2,-1.55) {乙：一条记录 \(=\) 一次会话};
  \foreach \x in {-4.8,-3.9,-3.0,-2.1,-1.4}{\fill[blue!65] (\x,-2.35) circle (0.16);}
  \foreach \x in {1.3,1.7,2.1,2.5,2.9,3.3,3.7,4.1,4.5,4.9,5.2}
    {\fill[blue!65] (\x,-2.35) circle (0.11);}
  \node[black!70] at (-3.2,-2.85) {\(\abs{\Omega}=5\)};
  \node[black!70] at (3.2,-2.85) {\(\abs{\Omega}=11\)};
\end{tikzpicture}

\emph{同一句话，两种``一条完整记录''，两个样本空间。后面的算术再无懈可击，也追不回这一步的分歧。}
\end{center}

这一单元把地基铺好，用的是概率建模的三件套 \((\Omega,\mathcal F,P)\)：先说清一次试验的结果有哪些，再说清哪些集合算得上事件，最后给这些事件指派满足可加性的数。三件套听起来像形式主义的排场，实际用处非常具体——三门问题、检验悖论、幸存者偏差、A/B 实验里对不上的两组数字，症结几乎都不在算错，而在两个人写下的 \(\Omega\) 根本不一样。把这句判断带在身上，比记住任何一条组合公式都管用。

第一篇\hyperref[sec:05-Probability/01-Probability-Spaces/01]{``随机实验、样本空间与事件''}处理前两件：一条完整记录该记到什么粒度，事件为什么是子集，以及日志粒度怎样提前决定了哪些问题日后还能问。第二篇\hyperref[sec:05-Probability/01-Probability-Spaces/02]{``概率公理与概率空间''}装上第三件，说明只用三条公理就能长出全部常用规则，也说明为什么可加性必须要到可数这一步。第三篇\hyperref[sec:05-Probability/01-Probability-Spaces/03]{``从现实问题到概率模型''}是这条主线的落地：从一段模糊的现实描述走到一个确定的概率空间，中途做了哪些取舍，等可能、独立、平稳这三条假设是怎么被悄悄塞进去的，以及它们失效时结论朝哪个方向坏。

三篇按依赖顺序排列，但从困惑进入同样可行。若眼下的麻烦是题面读不明白，从第一篇开始；若想知道概率规则的出处，直接看第二篇；若计算步骤都会做、答案却让人不安，第三篇最对症。

读完这一单元，一个合格的自检是能徒手写出一张模型卡片：一次完整结果长什么样，哪些事件可观察，概率由什么机制分配，以及改动粒度或观察方式之后，所问的还是不是同一个问题。写得出来，就可以进入\hyperref[ch:038]{``条件概率、Bayes 与独立性：新信息怎样改变判断''}——那里要处理的是新信息到来时，怎样把已经建好的样本空间重新裁剪一遍。
